% Diego - Ejercicio 3 inciso d.

\subsection{Inciso d. Inserciones con cardinalidad $1:1$}

Se considera el mismo escenario del inciso~b para las relaciones
\emph{A} y \emph{B}. Se toma como base el modelo relacional obtenido
para la cardinalidad $1:1$, suponiendo participación parcial de ambos
lados.

Las relaciones son
$\mathrm{A}(a_1,a_2,a_3)$ y $\mathrm{B}(b,b_1)$, mientras que la
relación $\mathrm{AB}(a_1,a_2,b,ab_1)$ tiene como llave primaria
$(a_1,a_2)$ y una restricción de unicidad sobre $b$.

En el escenario dado, las llaves de \emph{A} son
$(2,\text{'ww'})$, $(4,\text{'xx'})$, $(6,\text{'yy'})$ y
$(8,\text{'zz'})$, mientras que las llaves de \emph{B} son
$17$, $27$, $37$, $47$ y $57$.

\paragraph{Conjunto de tuplas insertables.}

Con los datos disponibles, no es posible proponer un conjunto de cinco
tuplas que pueda insertarse por completo en \emph{AB}. Esto se debe
a que existen únicamente cuatro llaves en \emph{A} y la llave primaria
$(a_1,a_2)$ de \emph{AB} impide que una misma entidad de \emph{A}
aparezca en más de una tupla de dicha relación.

El conjunto máximo de tuplas que puede insertarse simultáneamente
es el siguiente, expresado en el orden
$\mathrm{AB}(a_1,a_2,b,ab_1)$:

\begin{enumerate}[label=\roman*.,leftmargin=*,itemsep=0.2em]
    \item \texttt{(2,'ww',17,5)}
    \item \texttt{(4,'xx',27,10)}
    \item \texttt{(6,'yy',37,15)}
    \item \texttt{(8,'zz',47,20)}
\end{enumerate}

Las cuatro tuplas pueden insertarse porque cada pareja $(a_1,a_2)$
existe en \emph{A}, cada valor de $b$ existe en \emph{B}, los valores
de $ab_1$ son enteros y ninguna llave primaria ni valor sujeto a la
restricción de unicidad se repite.

La participación parcial de ambos lados permite que la entidad de
\emph{B} con llave $57$ permanezca sin asociarse, por lo que no es
obligatorio que todas las entidades de ambas relaciones participen
en \emph{AB}.

\paragraph{Conjunto de tuplas no insertables.}

Se propone el siguiente conjunto de cinco tuplas, también en el orden
$\mathrm{AB}(a_1,a_2,b,ab_1)$:

\begin{enumerate}[label=\roman*.,leftmargin=*,itemsep=0.2em]
    \item \texttt{(2,'ww',17,5)}
    \item \texttt{(4,'xx',27,10)}
    \item \texttt{(6,'yy',37,15)}
    \item \texttt{(8,'zz',47,20)}
    \item \texttt{(2,'ww',57,25)}
\end{enumerate}

Este conjunto no puede insertarse por completo. La quinta tupla
repite la llave primaria $(2,\text{'ww'})$ de la primera tupla.
Aunque el valor $57$ existe en \emph{B} y los valores de $ab_1$ son
enteros,la cardinalidad $1:1$ exige que una entidad de \emph{A} no
esté asociada con más de una entidad de \emph{B} a través de
\emph{AB}. Por lo tanto, la quinta tupla no puede insertarse.

\paragraph{Conclusión.}

La cardinalidad $1:1$ se garantiza mediante la llave primaria
$(a_1,a_2)$ y la restricción $\operatorname{UNIQUE}(b)$ en
\emph{AB}. La participación parcial permite que algunas entidades
no tengan asociación, pero no permite que una misma entidad tenga
más de una asociación. Debido a que \emph{A} contiene únicamente
cuatro entidades en el escenario planteado, el conjunto máximo
insertable tiene cuatro tuplas y no cinco.
